\chapter{Events you cannot see one by one}
\label{ch:problem}

\begin{motivation}
You manage a hospital and want to know whether infections are
\emph{self-propagating} (one case making the next more likely) or merely tracking
an external driver. You have a clean record of \emph{admissions per day} --- but
not who infected whom, nor the exact infection times. The model you want (a Hawkes
process) is written in terms of event \emph{times}; your data are \emph{counts}.
This chapter is about that mismatch and how we will dissolve it.
\end{motivation}

\section{Three problems, one shape}

Consider:
\begin{itemize}[leftmargin=1.4em]
\item \textbf{Epidemics.} Daily new cases or hospital admissions; individual
infection times are unknown (privacy, or simply unrecorded).
\item \textbf{Online popularity.} Daily views of a video, driven exogenously by
shares/tweets and endogenously by word-of-mouth; platforms report daily volumes,
not per-view timestamps.
\item \textbf{Markets / operations.} Trades or tickets aggregated per minute or
hour for storage or reporting.
\end{itemize}
In each case the data are \emph{interval-censored}: time is partitioned into
buckets $(o_{i-1},o_i]$ and we observe only the count
$\sfont{C}(o_{i-1},o_i]$ in each. We would like to fit a self-exciting model and
ask the same questions we could ask from timestamps: how much of the activity is
self-generated, how fast does the excitation decay, how do covariates modulate it?

\section{The model we want, and why it resists}

The Hawkes process has conditional intensity (rate of events given the past)
\begin{equation}
\lambda(t)=\underbrace{s(t)}_{\text{exogenous}}+\underbrace{\sum_{t_i<t}\phi(t-t_i)}_{\text{endogenous}},
\label{eq:ch1-hawkes}
\end{equation}
where $s\ge0$ injects \emph{immigrant} events and the kernel $\phi\ge0$ says each
past event lifts the future rate, creating \emph{offspring}. Two obstructions stop
us fitting \eqref{eq:ch1-hawkes} to counts.

\begin{whybox}[: timestamps]
The maximum-likelihood objective for a point process,
$\ell(\theta)=\sum_{t_i\le T}\log\lambda(t_i)-\int_0^T\lambda$, evaluates the
intensity \emph{at the event times} $t_i$. No times, no likelihood.
\end{whybox}

\begin{whybox}[: dependent increments]
Even if we tried to model just the counts, the Hawkes process has \emph{no
independent-increments property}: the number of events in a bucket depends on the
counts in earlier buckets (past events are still exciting the present). So we
cannot treat the buckets as independent Poisson variables either --- the trick that
makes ordinary count models easy is unavailable.
\end{whybox}

\section{The plan: fit a well-behaved twin}

We will not fit the Hawkes process directly. Instead we build a \emph{companion}
process --- the Mean Behavior Poisson process --- whose intensity is the
\emph{expected} Hawkes intensity,
\begin{equation}
\xi(t)=\E[\lambda(t)],
\label{eq:ch1-mbpp}
\end{equation}
and which, being an (inhomogeneous) Poisson process, \emph{does} have independent
increments and therefore a tractable count likelihood. Because $\xi$ is built from
the same $s$ and $\phi$ as the Hawkes process, fitting the companion to counts
recovers (approximations of) the Hawkes parameters.

\begin{intuitionbox}
Stop trying to track each random realisation and track its \emph{average} instead.
The average intensity is deterministic; a process with a deterministic rate is
Poisson; and Poisson counts are independent and easy. We trade a little fidelity
(the average is not the process) for a likelihood we can actually write down.
\end{intuitionbox}

\section{What this book builds}

By the end you will be able to: simulate self-exciting data with or without
covariates (Ch.~\ref{ch:hawkes}, \ref{ch:excitation}); compute the MBPP intensity
and compensator in closed form, by ODE, or by quadrature (Ch.~\ref{ch:solving},
\ref{ch:solvers}); fit the model from counts and quantify uncertainty
(Ch.~\ref{ch:fitting}, \ref{ch:identif}); let covariates drive the baseline and
even the excitation (Ch.~\ref{ch:baseline}, \ref{ch:excitation}); and replace the
solver with a learned neural operator when no closed form exists
(Ch.~\ref{ch:neural}). Crucially, you will know \emph{when each tool is the right
one} and \emph{what each one cannot tell you}.

\begin{recap}
Interval-censored counts break the Hawkes likelihood (no timestamps) and the
Poisson shortcut (no independent increments). The fix is to fit a deterministic-mean
companion process that restores both --- the subject of Part II. The two ingredients
we need first are the Poisson process (Ch.~\ref{ch:poisson}) and the Hawkes process
(Ch.~\ref{ch:hawkes}).
\end{recap}

\exercise{Name a dataset from your own field that is interval-censored. What would
the ``immigrants'' and ``offspring'' be, and why might the timestamps be
unavailable?}
\exercise{Explain in one sentence each, to a colleague, the two reasons the Hawkes
likelihood cannot be used on counts.}
